\section{Week 8: Human-Centred Values -- Environmental Sustainability}

\subsection{Topic Summary}

\begin{keybox}[Learning Objectives]
By the end of this week, you will be able to:
\begin{itemize}
  \item \textbf{HAI3}: Apply critical reasoning skills to identify general (explainability, fairness, efficacy, accountability, transparency) and context-specific constraints on human-AI interactions.
\end{itemize}
\end{keybox}

\begin{keybox}[Big Picture]
AI systems have significant environmental costs---energy, water, minerals, carbon emissions. This week explores:
\begin{enumerate}
  \item \textbf{The scale of the problem}: Data centers, training costs, inference costs.
  \item \textbf{Calculating AI energy usage}: Factors that affect environmental impact.
  \item \textbf{Eco-feedback tools}: For practitioners and end-users.
  \item \textbf{Practitioner perceptions}: Why action is difficult.
\end{enumerate}

The ``Green AI'' movement aims to estimate, analyze, and mitigate AI's environmental impact.
\end{keybox}

\begin{keybox}[What Examiners Like to Ask]
\begin{itemize}
  \item Explain the factors that affect AI's environmental footprint.
  \item Compare training vs.\ inference energy consumption.
  \item Describe how region and time affect carbon intensity.
  \item Give examples of eco-feedback tools for AI practitioners.
  \item Discuss barriers to sustainable AI practices (from practitioner research).
  \item Propose design interventions to make AI more sustainable.
\end{itemize}
\end{keybox}

\subsection{Overview + Core Concepts + Extended Theory}

\subsubsection*{The Scale of the Problem}

\paragraph*{Key Statistics}
\begin{itemize}
  \item \textbf{IEA 2024}: Data centers' electricity consumption $\geq$ 1000 TWh by 2026 (double from 2022).
  \item \textbf{IEA 2025}: Global data center consumption could reach 3500 TWh in 3 years---equivalent to adding Japan's electricity consumption annually.
  \item \textbf{US by 2028}: Data centers could use 12\% of total US electricity (currently $\approx$4.4\%).
\end{itemize}

\paragraph*{Real-World Impacts}
\begin{itemize}
  \item Delayed retirement of coal plants to meet demand.
  \item Conflicts with local communities over water and energy.
  \item West London: Data centers imbalanced grid; new buildings can't be powered.
  \item Water scarcity in Arizona, Spain, Ireland due to data center cooling.
\end{itemize}

\subsubsection*{Calculating AI Energy Usage}

\begin{defbox}[Green AI]
A movement to \textbf{estimate, analyze, and mitigate} the environmental impact of AI systems---including energy consumption, water usage, mineral extraction, and carbon emissions.
\end{defbox}

\paragraph*{Factors Affecting Environmental Impact}

\begin{center}
\renewcommand{\arraystretch}{1.3}
\begin{tabular}{@{} l p{8cm} @{}}
\toprule
\textbf{Factor} & \textbf{Impact} \\
\midrule
Training vs.\ Inference & Training is more intensive but one-time; inference adds up across billions of queries. \\
Cloud vs.\ Local & Cloud allows regional flexibility and batch scheduling. \\
Region & Energy grid mix varies (nuclear France vs.\ fossil-fuel US). \\
Time of day/season & Renewable availability varies (solar in summer, wind overnight). \\
Model architecture & Larger models = more computation = more energy. \\
Task type & Generative tasks (summarization, image generation) are far more intensive than discriminative (classification). \\
\bottomrule
\end{tabular}
\end{center}

\paragraph*{Life Cycle Analysis}

Three main stages contribute to environmental impact:

\begin{enumerate}
  \item \textbf{Equipment manufacturing}: Embodied emissions in hardware (servers, GPUs, chips).
  \item \textbf{Model training}: Most intensive phase; one-time but huge.
  \item \textbf{Model deployment/inference}: Lower per-query but scales with usage.
\end{enumerate}

\begin{exbox}[BLOOM Life Cycle Analysis (Luccioni et al., 2022)]
\textbf{BLOOM}: 176B parameter language model.

\textbf{Training emissions}:
\begin{itemize}
  \item Dynamic power only: $\approx$24.7 tons CO$_2$ equivalent.
  \item Full life cycle (equipment manufacturing + operational): $\approx$50 tons CO$_2$ equivalent.
\end{itemize}

\textbf{Lesson}: Training large models has substantial carbon footprint---equivalent to $\approx$5 round-trip flights NYC--London per ton.
\end{exbox}

\subsubsection*{Region and Time Matter}

\paragraph*{Energy Grid Mix}
Different regions have different carbon intensities:
\begin{itemize}
  \item \textbf{France}: Mostly nuclear $\to$ low carbon.
  \item \textbf{Germany/US/Australia}: More fossil fuels $\to$ high carbon.
\end{itemize}

\paragraph*{Time of Day/Season}
Renewable energy availability varies:
\begin{itemize}
  \item Solar: Higher during day, summer.
  \item Wind: Often higher at night, varies by season.
\end{itemize}

\textbf{Tool}: \texttt{electricitymaps.com} shows real-time carbon intensity by region.

\begin{tipbox}[Cloud Flexibility]
Cloud training allows you to:
\begin{itemize}
  \item Choose low-carbon regions.
  \item Schedule training during high-renewable periods.
  \item Batch tasks for efficiency.
\end{itemize}
\end{tipbox}

\subsubsection*{Task Type Matters}

\textbf{Generative tasks} are significantly more energy-intensive than \textbf{discriminative tasks}.

\begin{center}
\renewcommand{\arraystretch}{1.3}
\begin{tabular}{@{} l l @{}}
\toprule
\textbf{Lower Energy (Discriminative)} & \textbf{Higher Energy (Generative)} \\
\midrule
Text classification & Summarization \\
Token classification & Text generation \\
Image classification & Image generation \\
\bottomrule
\end{tabular}
\end{center}

\textbf{Note}: Difference is \textit{orders of magnitude} (logarithmic scale in original research).

\subsubsection*{Eco-Feedback for AI Practitioners}

\begin{defbox}[Eco-Feedback]
A device or visualization that provides feedback on environmental impact (e.g., energy consumption) to encourage more sustainable behavior.

\textbf{Hypothesis}: With more information, people will make more sustainable choices.
\end{defbox}

\paragraph*{Tools for Practitioners}

Research found 16+ tools that calculate, estimate, track, schedule, or label AI energy consumption:

\begin{center}
\renewcommand{\arraystretch}{1.3}
\begin{tabular}{@{} l p{8cm} @{}}
\toprule
\textbf{Tool} & \textbf{Description} \\
\midrule
CodeCarbon & Python package; tracks CO$_2$ emissions; dashboard over time; equivalents (households, miles driven). \\
EnergyVis & Interactive tracking; state-by-state US carbon intensity; compare hardware/locations. \\
AI Energy Score & HuggingFace leaderboard; compare foundational models by energy efficiency per task. \\
EcoLogits & Calculator + Python library; estimates energy, GHG, resources per prompt. \\
Chat UI Energy & Measures actual GPU energy during inference (open-source models only). \\
\bottomrule
\end{tabular}
\end{center}

\paragraph*{Industry Transparency}

\textbf{Mistral AI (2024)}: Published peer-reviewed life cycle assessment.
\begin{itemize}
  \item 400-token response: 1.14g CO$_2$ equivalent, 45ml water.
  \item Most emissions from training + inference; most materials from embodied impact.
\end{itemize}

\textbf{Google Gemini (2024)}: Published inference measurements.
\begin{itemize}
  \item Median Gemini prompt: 4,167--10,000 prompts per kWh (depends on metrics).
  \item Methodology: Active accelerator + CPU + idle time.
\end{itemize}

\begin{tipbox}[Measurements vs.\ Estimates]
Many figures are \textbf{estimates}, not measurements. When comparing:
\begin{itemize}
  \item Check if data is measured or estimated.
  \item Open-source models allow measurement; closed-source requires estimation.
\end{itemize}
\end{tipbox}

\subsubsection*{Perceptions of AI Practitioners}

\textbf{Research study}: Interviews with 23 ML practitioners (academia + industry).

\paragraph*{Key Findings}

\begin{center}
\renewcommand{\arraystretch}{1.3}
\begin{tabular}{@{} l p{8cm} @{}}
\toprule
\textbf{Theme} & \textbf{Example Quote / Finding} \\
\midrule
Minimal perceived impact & ``ML affects maybe 2--3\%''---seen as small compared to other sectors. \\
ML's value outweighs cost & ``I don't care how environmentally friendly it is, I want to save lives'' (medical domain). \\
Individual powerlessness & ``As an individual, you can't really do much about it.'' \\
Responsibility deflection & ``My work is just a very small amount of this unsustainable system.'' \\
Uncertainty about measurement & Hard to know actual impact; many tools give estimates, not measurements. \\
\bottomrule
\end{tabular}
\end{center}

\paragraph*{Paths to Change}
\begin{itemize}
  \item Individuals have \textit{some} responsibility (awareness at minimum).
  \item \textbf{Regulation} can drive action.
  \item Need more \textbf{transparency and open data} on model impacts.
  \item Learn from Responsible AI discourse (fairness, accountability).
\end{itemize}

\subsubsection*{Eco-Feedback for End Users}

End users (not just developers) consume AI through inference. Should they know the cost?

\paragraph*{Tools for End Users}

\begin{center}
\renewcommand{\arraystretch}{1.3}
\begin{tabular}{@{} l p{8cm} @{}}
\toprule
\textbf{Tool} & \textbf{Description} \\
\midrule
EcoLogits Calculator & Web app: input provider, model, prompt size $\to$ get energy, GHG, resources estimate. \\
Chat UI Energy & Live energy measurement during chat (open-source models); shows \% of phone charge, latency. \\
\bottomrule
\end{tabular}
\end{center}

\textbf{Design opportunity}: Incorporate eco-feedback into user-facing AI products to increase transparency.

\subsubsection*{Key Takeaways}

\begin{enumerate}
  \item \textbf{AI has real environmental costs}: Energy, water, minerals, carbon.
  \item \textbf{Scale matters}: 1000+ TWh by 2026; 12\% of US electricity by 2028.
  \item \textbf{Multiple factors affect impact}: Training vs.\ inference, region, time, model size, task type.
  \item \textbf{Generative $\gg$ discriminative}: Image generation uses orders of magnitude more than classification.
  \item \textbf{Cloud enables optimization}: Choose low-carbon regions, schedule during high-renewable periods.
  \item \textbf{Tools exist}: CodeCarbon, EcoLogits, AI Energy Score---use them.
  \item \textbf{Practitioners feel powerless}: Need regulation, transparency, cultural shift.
  \item \textbf{End users can be informed}: Eco-feedback in user-facing products.
\end{enumerate}

\subsubsection*{Key Definitions Summary}

\begin{center}
\renewcommand{\arraystretch}{1.3}
\begin{tabular}{@{} l p{8cm} @{}}
\toprule
\textbf{Term} & \textbf{Definition} \\
\midrule
Green AI & Movement to estimate, analyze, mitigate AI's environmental impact \\
Eco-feedback & Visualization/feedback to encourage sustainable behavior \\
Carbon intensity & CO$_2$ emissions per unit of electricity (varies by region/time) \\
Energy grid mix & Proportion of energy from renewables vs.\ fossil fuels \\
Embodied emissions & Carbon footprint of manufacturing hardware \\
Life cycle analysis & Full assessment from manufacturing to disposal \\
Discriminative task & Classification, labeling (lower energy) \\
Generative task & Text/image generation (higher energy) \\
\bottomrule
\end{tabular}
\end{center}

\subsubsection*{Recommended Tools \& Resources}
\begin{itemize}
  \item \textbf{CodeCarbon}: \texttt{codecarbon.io}
  \item \textbf{EcoLogits}: \texttt{huggingface.co/spaces/genai-impact/ecologits-calculator}
  \item \textbf{AI Energy Score}: \texttt{huggingface.co/spaces/AIEnergyScore/Leaderboard}
  \item \textbf{Electricity Maps}: \texttt{electricitymaps.com}
  \item \textbf{Paper}: Luccioni et al., ``Estimating the Carbon Footprint of BLOOM''
  \item \textbf{Paper}: Luccioni et al., ``Power Hungry Processing''
\end{itemize}


